% Generated by `taskgen course examples` from factoring.ec_mod_n_add (seed 3).
% Do not edit: change examples.yaml and regenerate.
\begin{example}[необратимый знаменатель при удвоении]\label{ex:ecm-double-fail}
Кривая $y^2 = x^3 + 49x + 114$ рассматривается над $\mathbb{Z}/143\mathbb{Z}$. Вычислить $2P$ для $P = (2, 99)$ по обычным формулам с делением по модулю $143$. Если деление невозможно, найти нетривиальный делитель числа $143$.

\par\noindent\textit{Решение.} Точка лежит на кривой (проверка подстановкой). Удвоение: $\lambda = \dfrac{3x_1^2 + a}{2y_1}$, знаменатель $d = 2 \cdot 99 \bmod 143 = 55$, числитель $3 \cdot 2^2 + 49 \equiv 61$. Алгоритм Евклида для $(143, 55)$: остатки $143, 55, 33, 22, 11, 0$, т.\,е. $\gcd(55, 143) = 11 > 1$: $55$ необратим по модулю $143$, и найден делитель $11$, $143 = 11 \cdot 13$.\par Причина: по модулю $11$ имеем $P \equiv (2, 0)$ --- точка порядка $2$, и $2P = \mathcal{O}$; по модулю $13$ результат --- обычная точка $(5, 9)$.\[143 = 11 \cdot 13.\]
\end{example}
